% !TeX root = ../组合数学笔记.tex
% 由主文档通过 \input 载入；请编译主文档。
\setcounter{chapter}{6} % 保留原书章号。
\chapter{普通生成函数}
\label{ch:ogf}

生成函数把一个序列的全部信息存入各次幂的系数，使递推和组合分解转化为代数运算。本章使用形式幂级数，不把 $x$ 看成必须代入数值的实变量。因此运算的合法性来自逐系数的有限性，而非解析收敛。

\section{形式幂级数与系数提取}

\begin{definition}
对有理数序列 $(a_n)_{n\ge0}$，形式表达式 $A(x)=\sum_{n\ge0}a_nx^n$ 称为其\emph{普通生成函数（ordinary generating function）}，记 $[x^n]A(x)=a_n$。所有这类表达式组成 $\mathbb Q[[x]]$；相等定义为各系数分别相等，加法逐项进行，乘法定义为
\[
 [x^n](A(x)B(x))=\sum_{k=0}^n a_kb_{n-k}.
\]
\end{definition}

每个乘积系数只涉及有限项，故乘法良定义，并满足通常的结合、交换和分配律。若 $a_0\ne0$，方程 $A(x)B(x)=1$ 逐项给出 $b_0=1/a_0$ 与 $b_n=-a_0^{-1}\sum_{k=1}^n a_kb_{n-k}$，所以乘法逆唯一存在；若 $a_0=0$ 则不可能有逆。特别地，$1/(1-x)=\sum_{n\ge0}x^n$ 是形式恒等式。

当 $B(0)=0$ 时，复合 $A(B(x))=\sum_{k\ge0}a_kB(x)^k$ 也良定义，因为 $B(x)^k$ 的最低可能次数至少为 $k$。若内层常数项非零，则同一次幂可能收到无限项贡献，不能自动作形式复合。系数可按需要扩展到 $\mathbb R$ 或 $\mathbb C$，以处理特征根；组合计数序列本身仍取整数。

\section{加法、乘法、移位与标记}

若两类对象互斥且大小定义相同，其并集对应生成函数相加。若对象唯一分为一对有序组件，大小分别为 $k,n-k$，且两组件可独立选择，则整体数为 $\sum_{k=0}^n a_kb_{n-k}$，对应生成函数相乘。这一运算称为\emph{Cauchy 卷积（Cauchy product）}。若组件共享同一批不同标签，还需先分配标签，普通卷积便不够，见第\ref{ch:egf}章。

移位必须处理起点。例如 $\sum_{n\ge1}a_{n-1}x^n=xA(x)$，而 $\sum_{n\ge0}a_{n+1}x^n=(A(x)-a_0)/x$。形式微分定义为 $A'(x)=\sum_{n\ge1}na_nx^{n-1}$，故 $xA'(x)$ 对应在每个大小为 $n$ 的对象中标记一个位置；它满足乘法法则，可以逐系数验证。常用展开
\begin{equation}
  \frac{1}{(1-x)^k}=\sum_{n\ge0}\binom{n+k-1}{k-1}x^n\quad(k\ge1)
  \label{eq:negative-binomial}
\end{equation}
来自 $k$ 个几何级数相乘，系数恰是隔板法计数。

\section{从递推到有理生成函数}

对式\eqref{eq:tiling-recurrence}，令 $A(x)=\sum_{n\ge0}a_nx^n$，将 $n\ge2$ 的递推乘以 $x^n$ 后求和，得到
\[
 A(x)-1-x=x(A(x)-1)+x^2A(x),\qquad
 A(x)=\frac{1}{1-x-x^2}.
\]
这里每个等式均逐系数成立。分母分解为 $(1-\phi x)(1-\psi x)$，部分分式展开配合几何级数可恢复上一章的闭式。

\begin{proposition}
生成函数 $A(x)$ 可写成 $P(x)/Q(x)$，其中 $P,Q$ 为多项式且 $Q(0)=1$，当且仅当其系数从某一项起满足常系数线性递推。
\end{proposition}
\begin{proof}
写 $Q(x)=1-c_1x-\cdots-c_rx^r$，比较 $Q(x)A(x)=P(x)$ 中超过 $\deg P$ 且至少为 $r$ 的系数，得到递推。反之，将从某项起成立的递推移到一边，乘以 $x^n$ 求和；除有限个初始系数外均抵消，所以 $Q(x)A(x)$ 是多项式。
\end{proof}

若 $r$ 阶递推从 $n=r$ 起成立且 $c_r\ne0$，所得分子次数小于 $r$。在复数域分解 $Q(x)=\prod_j(1-\lambda_jx)^{m_j}$，其中不同 $\lambda_j$ 为非零特征根。部分分式没有多项式部分，而
\[
 [x^n](1-\lambda x)^{-h}=\binom{n+h-1}{h-1}\lambda^n.
\]
于是每个重数为 $m_j$ 的根贡献一个次数至多 $m_j-1$ 的多项式乘 $\lambda_j^n$。这证明了上一章给出的重根形式；分式恒等式的唯一性也保证这些项足以且唯一地表达由初值确定的序列。

\section{分配、零钱兑换与多变量生成函数}

变量 $x_i$ 允许取 $0,\ldots,u_i$ 时，其大小生成函数为 $1+x+\cdots+x^{u_i}$，所以总和为 $n$ 的解数是
\[
 [x^n]\prod_{i=1}^k(1+x+\cdots+x^{u_i})
 =[x^n]\frac{\prod_i(1-x^{u_i+1})}{(1-x)^k}.
\]
展开分子并用式\eqref{eq:negative-binomial}，正好得到容斥式\eqref{eq:bounded-compositions}。两种方法不同，但计数的是同一组整数向量。

\begin{example}[零钱兑换]
用面值 $1,2,5$ 的硬币兑换 $10$，每种不限数量，同面值硬币不可区分。生成函数为 $1/((1-x)(1-x^2)(1-x^5))$。固定五元硬币数为 $c=0,1,2$，二元硬币数分别有 $6,3,1$ 种，一元硬币数随即确定，故共有 $10$ 种。若把投币先后次序也视为方案的一部分，模型已改为字符串，不能沿用这个乘积。
\end{example}

再以 $y$ 记录硬币枚数，生成函数变为 $1/((1-yx)(1-yx^2)(1-yx^5))$，$[x^{10}y^4]$ 计数恰用四枚凑成 $10$ 的方案。每种面值为正数，保证固定总金额只涉及有限个配置。一般多变量生成函数应明确每个变量标记的统计量，提取系数后再将其他变量赋值。

\section{习题}
\begin{enumerate}
  \item\label{ex:ogf-1} 写出 $x_1+x_2+x_3=n$ 中 $x_1$ 为非负偶数、$x_2\ge1$、$0\le x_3\le2$ 的普通生成函数，并求 $n=5$ 的解数。
  \item\label{ex:ogf-2} 从 $a_n=2a_{n-1}+1,a_0=0$ 推导生成函数并提取系数。
  \item\label{ex:ogf-3} 用面值 $1,2,5$ 的硬币恰好四枚凑成 $10$，求方案数；写出相应双变量系数表达式。
  \item\label{ex:ogf-4} 形式等式 $1/(1-x)=\sum_{n\ge0}x^n$ 是否允许直接令 $x=2$ 并把右侧理解为通常的数项级数？解释原因。
\end{enumerate}

